\documentclass[11pt,letterpaper]{article}

\usepackage[margin=1in]{geometry}
\usepackage{graphicx}
\usepackage{booktabs}
\usepackage{amsmath}
\usepackage{amssymb}
\usepackage{hyperref}
\usepackage{xcolor}
\usepackage{caption}
\usepackage{float}
%\usepackage{siunitx}  % not available on this system
\providecommand{\num}[1]{#1}  % fallback for \num when siunitx is absent
%\usepackage{multirow}  % not available on this system
%\usepackage{enumitem}  % not available on this system

\hypersetup{
  colorlinks=true,
  linkcolor=blue,
  citecolor=blue,
  urlcolor=blue
}

% --- Local macros (subset of PPG12 defs.sty) ---
\newcommand{\pp}{\mbox{$p$+$p$}}
\newcommand{\sqs}{\mbox{$\sqrt{s}$}}
\newcommand{\pT}{\ensuremath{p_{\text{T}}}}
\newcommand{\ET}{\ensuremath{E_{\text{T}}}}
\newcommand{\zvtx}{\ensuremath{z_{\mathrm{vtx}}}}
\newcommand{\abseta}{\ensuremath{|\eta|}}
\newcommand{\isoETtruth}{\ensuremath{E_\text{T}^\text{iso,truth}}}
\newcommand{\Rcone}{\ensuremath{R_{\text{cone}}}}
\newcommand{\DR}{\ensuremath{\Delta R}}

\graphicspath{{../../plotting/figures/}}

\title{%
  \vspace{-1cm}
  {\large \textbf{sPHENIX Internal Note}}\\[0.5cm]
  {\LARGE \textbf{Eta Edge Migration Study:\\
  Fake Rate from Fiducial Boundary Effects}}\\[0.3cm]
  {\large Isolated Photon Cross-Section in \pp{} at $\sqs = 200$~GeV}
}
\author{sPHENIX Collaboration}
\date{\today}

\begin{document}

\maketitle

\begin{abstract}
This note documents a study of pseudorapidity edge migration effects in the
sPHENIX isolated photon cross-section measurement (PPG12). Both the truth and
reconstructed $\eta$ fiducial cuts are set at $\abseta < 0.7$, which allows
truth photons produced outside the fiducial volume to have reconstructed
clusters migrate inside, and vice versa. In the unfolding framework, inward-migrated
clusters matched to truth photons outside $\abseta < 0.7$ should be treated as
fakes but currently enter the response matrix as signal in
\texttt{RecoEffCalculator\_TTreeReader.C}. We quantify this contamination for
both single-interaction and double-interaction MC, finding a signal region (ABCD
region~A) contamination of ${\sim}1.2\%$ for single interaction and
${\sim}0.3\%$ for double interaction, with significantly larger effects in the
non-tight ABCD regions used for background estimation.
\end{abstract}

\tableofcontents
\newpage

%% ============================================================
\section{Introduction}
\label{sec:intro}

The PPG12 analysis measures the isolated prompt photon production cross-section
in \pp{} collisions at $\sqs = 200$~GeV using the sPHENIX EMCal. The fiducial
pseudorapidity acceptance is $\abseta < 0.7$, corresponding to the barrel EMCal
coverage. This cut is applied identically to both truth-level photons and
reconstructed clusters.

Because the $\eta$ boundary is a sharp geometric cut rather than a smooth
efficiency edge, photons near $|\eta| \approx 0.7$ can migrate across the
boundary between truth and reconstruction:
\begin{itemize}
  \item \textbf{Inward migration (fakes):} Truth photons with
    $|\eta^\text{truth}| > 0.7$ produce reconstructed clusters with
    $|\eta^\text{reco}| < 0.7$. These enter the signal sample but have no
    corresponding truth entry within the fiducial volume.
  \item \textbf{Outward migration (losses):} Truth photons with
    $|\eta^\text{truth}| < 0.7$ produce reconstructed clusters with
    $|\eta^\text{reco}| > 0.7$, escaping the reconstructed sample.
\end{itemize}

In the current implementation of \texttt{RecoEffCalculator\_TTreeReader.C},
inward-migrated clusters enter the response matrix as signal --- i.e., they are
treated as correctly matched truth-reco pairs. This is incorrect: their truth
$\pT$ falls outside the fiducial region, biasing the unfolding. This study
quantifies the magnitude of this effect.

%% ============================================================
\section{Method}
\label{sec:method}

A dedicated macro \texttt{EtaMigrationStudy.C} processes the full photon MC
samples: \texttt{photon5} (truth $\pT$ range $0$--$14$~GeV), \texttt{photon10}
($14$--$30$~GeV), and \texttt{photon20} ($30$+~GeV). Each sample is weighted by
its production cross-section.

Truth photons are registered with \emph{no} $\eta$ filter --- all
direct/fragmentation photons satisfying $\isoETtruth < 4$~GeV (cone
$\Rcone = 0.3$) are included regardless of $\eta$. Each truth--reco matched
pair is then classified into four categories based on the fiducial boundary
$\abseta < 0.7$:
\begin{enumerate}
  \item \textbf{Good:} both truth and reco inside $\abseta < 0.7$.
  \item \textbf{Fake (inward):} truth outside, reco inside.
  \item \textbf{Lost (outward):} truth inside, reco outside.
  \item \textbf{Both outside:} neither truth nor reco in fiducial volume.
\end{enumerate}

The standard PPG12 ABCD classification is applied to each pair:
\begin{itemize}
  \item Region A (tight + isolated): signal region.
  \item Region B (tight + non-isolated): background template.
  \item Region C (non-tight + isolated): background control.
  \item Region D (non-tight + non-isolated): background normalization.
\end{itemize}

For the double-interaction study, a toy vertex shift is applied: a random second
vertex is drawn from the data vertex distribution, and the new vertex is
computed as $z_\text{vtx}^\text{new} = (z_\text{vtx}^\text{orig} +
z_\text{vtx}^\text{rand}) / 2$. Cluster kinematics are recalculated using the
CEMC cylinder geometry ($R = 93.5$~cm):
\begin{equation}
  z/r_\text{new} = z/r_\text{old} + \frac{z_\text{vtx}^\text{old} - z_\text{vtx}^\text{new}}{93.5~\text{cm}},
  \quad
  \eta_\text{new} = \sinh^{-1}(z/r_\text{new}),
  \quad
  \ET^\text{new} = E / \cosh(\eta_\text{new}).
\end{equation}

%% ============================================================
\section{Results}
\label{sec:results}

%% --------------------------------------------------
\subsection{Overall Migration Rates}
\label{sec:overall}

Table~\ref{tab:overall} summarizes the unweighted migration rates per MC sample.
The fake (inward) rate is defined as the number of inward-migrated pairs divided
by the total number of matched pairs, and similarly for the loss (outward) rate.

\begin{table}[htbp]
\centering
\caption{Overall migration rates (unweighted) for single-interaction MC. Fake
  rate = inward / total; loss rate = outward / total.}
\label{tab:overall}
\begin{tabular}{lrrrrrcc}
\toprule
Sample & Total & Good & Fakes & Losses & Both out & Fake rate & Loss rate \\
\midrule
photon5  & \num{2361432} & \num{1523346} & \num{18490} & \num{13687} & \num{805909} & 1.20\% & 0.89\% \\
photon10 & \num{532878}  & \num{374380}  & \num{3942}  & \num{5822}  & \num{148734} & 1.04\% & 1.53\% \\
photon20 & \num{128595}  & \num{107741}  & \num{721}   & \num{2193}  & \num{17940}  & 0.66\% & 1.99\% \\
\bottomrule
\end{tabular}
\end{table}

The fake rate decreases with increasing truth $\pT$ range (from 1.20\% to
0.66\%), while the loss rate increases (from 0.89\% to 1.99\%). The asymmetry
arises because higher-$\pT$ photons have narrower shower profiles and more
precise $\eta$ reconstruction, but the steeply falling spectrum preferentially
populates the inward direction at lower $\pT$.

%% --------------------------------------------------
\subsection{Signal Region (A) Contamination}
\label{sec:signal_contamination}

The physically relevant metric is the contamination of the ABCD signal region
(region~A: tight BDT + isolated). This is defined as:
\begin{equation}
  \text{Contamination} = \frac{N_\text{fakes}^\text{tight+iso}}{N_\text{good}^\text{tight+iso} + N_\text{fakes}^\text{tight+iso}}.
  \label{eq:contamination}
\end{equation}
This quantity directly measures the fraction of response matrix entries that are
edge-migration fakes.

\subsubsection{Single Interaction}

Table~\ref{tab:signal_single} shows the signal region contamination per $\pT$
bin for single-interaction MC.

\begin{table}[htbp]
\centering
\caption{Signal region (tight + isolated) contamination from $\eta$ edge
  migration, single-interaction MC. Counts are cross-section weighted.}
\label{tab:signal_single}
\begin{tabular}{cccc}
\toprule
$\pT$ bin (GeV) & Good (tight+iso) & Fakes (tight+iso) & Contamination \\
\midrule
$8$--$10$   & \num{1.04e8} & \num{1.30e6} & 1.24\% \\
$10$--$12$  & \num{3.71e7} & \num{4.43e5} & 1.18\% \\
$12$--$14$  & \num{1.45e7} & \num{2.06e5} & 1.40\% \\
$14$--$16$  & \num{6.28e6} & \num{7.87e4} & 1.24\% \\
$16$--$18$  & \num{2.86e6} & \num{3.76e4} & 1.30\% \\
$18$--$20$  & \num{1.35e6} & \num{1.79e4} & 1.31\% \\
$20$--$22$  & \num{6.62e5} & \num{8.68e3} & 1.29\% \\
$22$--$24$  & \num{3.41e5} & \num{4.42e3} & 1.28\% \\
$24$--$26$  & \num{1.79e5} & \num{2.34e3} & 1.30\% \\
$26$--$28$  & \num{9.19e4} & \num{1.61e3} & 1.72\% \\
$28$--$32$  & \num{7.70e4} & \num{1.49e3} & 1.90\% \\
$32$--$36$  & \num{2.11e4} & \num{5.57e2} & 2.57\% \\
\midrule
\textbf{Overall} & \num{1.67e8} & \num{2.10e6} & \textbf{1.24\%} \\
\bottomrule
\end{tabular}
\end{table}

The contamination is approximately flat at ${\sim}1.2$--$1.3\%$ across most of
the $\pT$ range but rises to 2.57\% in the highest bin ($32$--$36$~GeV).

\subsubsection{Double Interaction}

Table~\ref{tab:signal_double} shows the equivalent for double-interaction
(toy vertex shift).

\begin{table}[htbp]
\centering
\caption{Signal region (tight + isolated) contamination from $\eta$ edge
  migration, double-interaction MC (toy vertex shift).}
\label{tab:signal_double}
\begin{tabular}{cccc}
\toprule
$\pT$ bin (GeV) & Good (tight+iso) & Fakes (tight+iso) & Contamination \\
\midrule
$8$--$10$   & \num{7.81e6} & \num{2.58e4} & 0.33\% \\
$10$--$12$  & \num{2.78e6} & \num{1.12e4} & 0.40\% \\
$12$--$14$  & \num{1.00e6} & \num{3.37e3} & 0.33\% \\
$14$--$16$  & \num{4.09e5} & \num{1.33e3} & 0.33\% \\
$16$--$18$  & \num{1.79e5} & \num{3.19e2} & 0.18\% \\
$18$--$20$  & \num{8.37e4} & \num{2.66e2} & 0.32\% \\
$20$--$22$  & \num{4.00e4} & \num{2.13e2} & 0.53\% \\
$22$--$24$  & \num{2.07e4} & \num{1.07e2} & 0.51\% \\
$24$--$26$  & \num{1.10e4} & 0             & 0.00\% \\
$26$--$28$  & \num{5.82e3} & 53            & 0.91\% \\
$28$--$32$  & \num{4.87e3} & 1             & 0.02\% \\
$32$--$36$  & \num{1.39e3} & 1             & 0.07\% \\
\midrule
\textbf{Overall} & \num{1.24e7} & \num{4.27e4} & \textbf{0.34\%} \\
\bottomrule
\end{tabular}
\end{table}

The double-interaction contamination in region~A is significantly \emph{lower}
than single interaction ($0.34\%$ vs.\ $1.24\%$). This counterintuitive result
is explained in Section~\ref{sec:double_enhancement}.

%% --------------------------------------------------
\subsection{ABCD Region Breakdown}
\label{sec:abcd}

The edge migration does not affect all ABCD regions equally.
The contamination varies strongly by ABCD region, particularly in the
highest $\pT$ bins where the effect is most pronounced.

Region~C (non-tight + isolated) exhibits dramatically higher contamination,
reaching up to 52\% at $32$--$36$~GeV for single interaction. This occurs
because edge-migrated clusters have distorted shower shapes: the $\eta$
mismatch between truth and reco positions distorts the transverse shower profile,
causing these clusters to preferentially fail the tight BDT selection and land in
the non-tight category.

The signal region~A contamination remains modest ($1.2$--$2.6\%$) precisely
because the BDT effectively rejects edge fakes --- they carry anomalous shower
shapes that the BDT was trained to discriminate against.

%% --------------------------------------------------
\subsection{Double Interaction Enhancement}
\label{sec:double_enhancement}

The overall \emph{raw} inward fake rate (before ABCD classification) is strongly
amplified by double interaction: $1.2\%$ (single) vs.\ $13.7\%$ (double),
approximately an $11\times$ amplification. The vertex displacement in
double-interaction events shifts the reconstructed $\eta$ of all clusters,
dramatically increasing the probability of boundary crossing.

However, the signal region contamination moves in the \emph{opposite} direction:
$1.24\%$ (single) $\to$ $0.34\%$ (double). The explanation is that
vertex-shifted clusters acquire additional shower shape distortion on top of the
edge migration effect. The compounded distortion causes the tight BDT to reject
these clusters even more aggressively, suppressing their entry into region~A.

Figure~\ref{fig:signal_leakage} quantifies this by showing the fraction of
ABCD-classified fakes (i.e., those passing common quality cuts) that leak into
regions~A and~C, as a function of $\ET$.
For single interaction, $\sim\!50\%$ of ABCD-classified fakes pass tight+iso
(region~A) and $\sim\!40\%$ land in region~C. For double interaction, a
\emph{higher} fraction ($\sim\!70$--$85\%$) passes tight+iso --- a seemingly
counterintuitive result. The explanation is a \emph{selection effect}: the
common quality cuts reject $98.4\%$ of double-interaction inward fakes (vs.\
$81\%$ for single), acting as a much stronger pre-filter. The $1.6\%$ of double
fakes that survive are the ``best'' fakes --- truth photons barely outside the
boundary with small vertex shifts --- and these naturally also tend to pass the
tight BDT. The common cuts are doing the heavy lifting; their survivors are
biased toward photon-like signatures. The absolute signal region contamination
(Fig.~\ref{fig:abcd} vs.\ Fig.~\ref{fig:abcd_double}) remains much lower for
double interaction ($0.34\%$ vs.\ $1.24\%$) because the total number of
ABCD-classified fakes is ${\sim}60\times$ smaller.

It should be noted that the double-interaction statistics are marginal: only 77
unweighted tight+iso fake entries contribute to the entire double-interaction
signal region, limiting the precision of the contamination estimate.

%% --------------------------------------------------
\subsection{Response Matrix Contamination}
\label{sec:response}

The response matrix in \texttt{RecoEffCalculator\_TTreeReader.C} is filled for
all tight+iso matched pairs. The $1.2\%$ signal region contamination from
Table~\ref{tab:signal_single} directly translates to $1.2\%$ of response matrix
entries having incorrect truth $\pT$ (from outside the fiducial $\eta$ region).
These entries map reco $\pT$ values to truth $\pT$ values that should not
participate in the unfolding, introducing a bias in the unfolded cross-section.

The effect is $\pT$-dependent: rising from ${\sim}1.2\%$ at low $\pT$ to
${\sim}2.6\%$ at $32$--$36$~GeV. This $\pT$ dependence means the bias is not a
simple overall normalization shift but distorts the shape of the unfolded
spectrum.

%% --------------------------------------------------
\subsection{ABCD Leakage Fractions: Single vs.\ Double Interaction}
\label{sec:leakage_fractions}

The ABCD signal leakage fractions $c_B$, $c_C$, $c_D$ (defined as the ratio of
truth-matched signal in regions B, C, D to that in region~A) directly enter the
background subtraction and purity correction.  Figure~\ref{fig:leakage_fractions}
shows these fractions computed from the ``good'' signal (truth inside fiducial
$\abseta < 0.7$) for both single and double interaction.

\begin{itemize}
  \item \textbf{$c_B$ (tight, non-isolated):}  Increases from ${\sim}6$--$11\%$
    (single) to ${\sim}11$--$29\%$ (double), a factor ${\sim}2$--$3\times$.
    The vertex shift in double-interaction events smears the topo-cluster
    isolation sum, causing isolated photons to fail the isolation cut and migrate
    from region~A to B.

  \item \textbf{$c_C$ (non-tight, isolated):}  Comparable between single and
    double (${\sim}2$--$8\%$), with no systematic amplification.  The shower
    shape BDT is less sensitive to vertex shifts at the level relevant here.

  \item \textbf{$c_D$ (non-tight, non-isolated):}  Small in both cases
    ($< 2\%$), with modest increase for double interaction.
\end{itemize}

The dominant effect is the $c_B$ amplification.  Since the ABCD background
estimate is $N_\text{bkg}^A = R \cdot B \cdot C / D$ with signal leakage
corrections proportional to $(1 - c_B)(1 - c_C)/(1 - c_D)$, the $2$--$3\times$
increase in $c_B$ for double interaction significantly affects the corrected
purity, particularly at high $\ET$ where the leakage is largest.

%% ============================================================
\section{Discussion}
\label{sec:discussion}

The ${\sim}1.2\%$ signal region contamination is a systematic effect, not a
statistical fluctuation. At the highest $\pT$ bin ($32$--$36$~GeV), the
contamination reaches 2.6\%, which is comparable in magnitude to other
systematic uncertainties in the analysis (e.g., energy scale, BDT threshold
variations).

Three approaches could address this:

\begin{enumerate}
  \item \textbf{Asymmetric $\eta$ cuts:} Apply a tighter reco $\eta$ cut than
    truth (e.g., $|\eta^\text{reco}| < 0.6$ with $|\eta^\text{truth}| < 0.7$).
    This eliminates inward migration but sacrifices acceptance.

  \item \textbf{Truth $\eta$ check in response matrix:} Explicitly verify that
    the truth $\eta$ of matched pairs satisfies $|\eta^\text{truth}| < 0.7$
    before filling the response matrix. Outside-fiducial matches are counted as
    fakes rather than signal. This is the minimal fix.

  \item \textbf{Systematic uncertainty:} Include the contamination as a
    systematic uncertainty on the unfolded cross-section, propagated through the
    full pipeline.
\end{enumerate}

The region~C contamination (up to 52\% at high $\pT$) is a separate concern.
Region~C enters the ABCD background estimation through the $R$-factor:
$R = N_C / N_D$. If a large fraction of region~C entries are edge-migration
fakes rather than genuine non-tight background, the $R$-factor and consequently
the background subtraction could be biased.

For double interaction, the tight+iso selection provides effective rejection of
edge-migrated clusters (contamination drops from $13.7\%$ raw to $0.34\%$ in
region~A). However, the ${\sim}14\%$ raw migration rate is large enough to
significantly affect the non-tight regions used in background estimation.

%% ============================================================
\section{Recommendations}
\label{sec:recommendations}

\begin{enumerate}
  \item \textbf{Short term:} Add the truth $\eta$ check to the response matrix
    filling in \texttt{RecoEffCalculator\_TTreeReader.C}. Matched pairs with
    $|\eta^\text{truth}| > 0.7$ should be counted as fakes, not signal.

  \item \textbf{Medium term:} Evaluate the impact on the unfolded cross-section
    by running the full pipeline with and without the fix. The difference
    quantifies the bias introduced by the current implementation.

  \item \textbf{Background estimation:} Evaluate and include the region~C
    contamination as a systematic uncertainty on the ABCD $R$-factor.
\end{enumerate}

%% ============================================================
\section{Appendix: Plots}
\label{sec:plots}

Figures~\ref{fig:matrix}--\ref{fig:deta} show the distributions produced by
\texttt{plot\_eta\_migration.C}.

\begin{figure}[htbp]
  \centering
  \begin{minipage}[t]{0.48\textwidth}
    \centering
    \includegraphics[width=\textwidth]{eta_migration_matrix.pdf}
    \caption{Migration matrix: truth vs.\ reco $\eta$ classification for
      single-interaction MC.}
    \label{fig:matrix}
  \end{minipage}
  \hfill
  \begin{minipage}[t]{0.48\textwidth}
    \centering
    \includegraphics[width=\textwidth]{eta_migration_matrix_double.pdf}
    \caption{Migration matrix: truth vs.\ reco $\eta$ classification for
      double-interaction MC (toy vertex shift).}
    \label{fig:matrix_double}
  \end{minipage}
\end{figure}

\begin{figure}[htbp]
  \centering
  \begin{minipage}[t]{0.48\textwidth}
    \centering
    \includegraphics[width=\textwidth]{eta_migration_fake_rate.pdf}
    \caption{Inward fake rate vs.\ $\ET$ for single and double interaction.}
    \label{fig:fake_rate}
  \end{minipage}
  \hfill
  \begin{minipage}[t]{0.48\textwidth}
    \centering
    \includegraphics[width=\textwidth]{eta_migration_loss_rate.pdf}
    \caption{Outward loss rate vs.\ $\ET$ for single and double interaction.}
    \label{fig:loss_rate}
  \end{minipage}
\end{figure}

\begin{figure}[htbp]
  \centering
  \begin{minipage}[t]{0.48\textwidth}
    \centering
    \includegraphics[width=\textwidth]{eta_migration_response_contamination.pdf}
    \caption{All-cluster inward migration fraction vs.\ $\ET$ for single
      (${\sim}2$--$8\%$) and double (${\sim}25\%$) interaction.
      This is the raw rate before any ABCD selection.}
    \label{fig:response}
  \end{minipage}
  \hfill
  \begin{minipage}[t]{0.48\textwidth}
    \centering
    \includegraphics[width=\textwidth]{eta_migration_tightiso_response.pdf}
    \caption{Response matrix impact: tight+iso inward contamination
      (filled circles) and outward loss rate (open squares) vs.\ $\ET$.
      Single interaction: ${\sim}1.2\%$ contamination and ${\sim}1.5\%$ loss.
      Double: ${\sim}0.3\%$ contamination but ${\sim}13\%$ loss.}
    \label{fig:tightiso_response}
  \end{minipage}
\end{figure}

\begin{figure}[htbp]
  \centering
  \begin{minipage}[t]{0.48\textwidth}
    \centering
    \includegraphics[width=\textwidth]{eta_migration_abcd_fakes.pdf}
    \caption{ABCD distribution of inward-migration fakes for single interaction.
      Region~A (tight+iso) and C (nontight+iso) dominate.}
    \label{fig:abcd}
  \end{minipage}
  \hfill
  \begin{minipage}[t]{0.48\textwidth}
    \centering
    \includegraphics[width=\textwidth]{eta_migration_abcd_fakes_double.pdf}
    \caption{ABCD distribution of inward-migration fakes for double interaction.
      The total yield is ${\sim}60\times$ smaller than single interaction
      (Fig.~\ref{fig:abcd}), because 98\% of double-interaction fakes fail
      common quality cuts.}
    \label{fig:abcd_double}
  \end{minipage}
\end{figure}

\begin{figure}[htbp]
  \centering
  \begin{minipage}[t]{0.48\textwidth}
    \centering
    \includegraphics[width=\textwidth]{eta_migration_signal_leakage.pdf}
    \caption{Fraction of inward-migrating fakes that leak into regions~A
      (tight+iso, filled circles) and C (non-tight+iso, open squares) as a
      function of $\ET$, comparing single (black) and double (red) interaction.
      Single-interaction fakes have a higher probability of passing tight+iso,
      while double-interaction fakes are more efficiently rejected into region~C.}
    \label{fig:signal_leakage}
  \end{minipage}
  \hfill
  \begin{minipage}[t]{0.48\textwidth}
    \centering
    \includegraphics[width=\textwidth]{eta_migration_truth_eta_fakes.pdf}
    \caption{Truth $\eta$ distribution of inward-migrated fakes.}
    \label{fig:truth_eta}
  \end{minipage}
\end{figure}

\begin{figure}[htbp]
  \centering
  \includegraphics[width=0.48\textwidth]{eta_migration_deta_vs_ET.pdf}
  \caption{$\Delta\eta$ (truth $-$ reco) vs.\ $\ET$ for edge-migrated
    clusters.}
  \label{fig:deta}
\end{figure}

\begin{figure}[htbp]
  \centering
  \includegraphics[width=0.55\textwidth]{eta_migration_leakage_fractions.pdf}
  \caption{ABCD signal leakage fractions $c_X = N_X^\text{sig} / N_A^\text{sig}$
    for single interaction (filled markers) and double interaction (open markers)
    as a function of $\ET$.  Upper panel: absolute leakage fractions for regions
    B, C, and D.  Lower panel: double-to-single ratio.  The isolation leakage
    $c_B$ increases by a factor of ${\sim}2$--$3\times$ under double interaction
    due to vertex-shift smearing of the isolation cone, while the shower-shape
    leakage $c_C$ is comparable or slightly reduced.}
  \label{fig:leakage_fractions}
\end{figure}

%% ============================================================
\clearpage
\section{Sideband Flow from Background-Only MC}
\label{sec:bkg_flow}

The results so far quantify the impact of eta edge migration on the signal
sample using photon MC. The PPG12 ABCD background estimate relies on three
sideband regions (B, C, D) that are populated overwhelmingly by jet-induced
fake clusters in data. To assess how fiducial boundary migration distorts
these sideband populations --- and therefore the background estimate itself
--- we run \texttt{EtaMigrationStudy.C} on the background-only MC with
\emph{every} truth particle registered as a candidate match (not just
direct/fragmentation photons), using \texttt{cluster\_truthtrkID} to find the
dominant truth contributor per cluster --- typically a decay photon from a
$\pi^{0}$ or $\eta$ in jet MC.

Two interaction scenarios are studied:
\begin{description}
  \item[Single-interaction background] jet20 ($21 < \pT^\text{jet} < 32$) and
    jet30 ($32 < \pT^\text{jet} < 42$) samples, cross-section weighted and
    merged. This matches the nominal PPG12 background MC.
  \item[Double-interaction background] the full-GEANT pileup sample
    \texttt{jet12\_double} ($\pT^\text{jet} > 10$), simulated with two
    overlapping \pp{} collisions at nominal luminosity. This represents the
    asymptotic ``100\% pileup'' rate; in real data at 1.5~mrad crossing only
    ${\sim}7.9\%$ of clusters come from double interactions.
\end{description}

\subsection{The NPB-Sentinel Issue}
\label{subsec:bkg_npb_sentinel}

The PPG12 ABCD common preselection requires an NPB (non-photon-background) TMVA
score $> 0.5$. The NPB model was trained only on clusters with $\abseta < 0.7$;
for clusters reconstructed outside the fiducial, the stored NPB score is the
sentinel value $-1$. Consequently the NPB cut rejects \emph{every} cluster
outside the fiducial volume, which eliminates outward-migration candidates from
the shower-shape sidebands before any flow measurement can be made. To restore
the symmetry, the eta-migration analysis fills a parallel ``no-NPB'' histogram
set in which the NPB common cut is skipped. The inflow (measured) side is
essentially unchanged because NPB rejects almost no in-fiducial background
clusters (the score is typically $> 0.9$ for real showers). The outflow side
now carries meaningful statistics.

All plots in this section use the no-NPB classification.

\subsection{Definitions}
\label{subsec:bkg_defs}

For each cluster with a valid truth match, we classify into one of three flow
categories per ABCD region $X \in \{A, B, C, D\}$:
\begin{description}
  \item[Native] truth and reco both inside $\abseta < 0.7$, cluster passes $X$.
    These are the clusters that \emph{should} count in region $X$.
  \item[Inflow] truth $\abseta > 0.7$, reco $\abseta < 0.7$, cluster passes $X$.
    Spurious additions to measured $X$ --- migration \emph{into} the fiducial
    sideband.
  \item[Outflow] truth $\abseta < 0.7$, reco $\abseta > 0.7$, cluster passes
    $X$ (by shower-shape/isolation criteria, independent of reco $\eta$).
    Clusters that would have populated region $X$ had the reco $\eta$ been
    correctly inside fiducial --- migration \emph{out of} the fiducial
    sideband. Note that the analysis does \emph{not} count these in the
    measured sideband, because the fiducial cut is applied on reco $\eta$.
\end{description}

We then define two flow rates:
\begin{align}
  r_\text{in}^{X}(\ET) &= \frac{N_X^\text{inflow}}
                              {N_X^\text{native} + N_X^\text{inflow}},
  \label{eq:bkg_inflow_rate} \\
  r_\text{out}^{X}(\ET) &= \frac{N_X^\text{outflow}}
                              {N_X^\text{native} + N_X^\text{outflow}}.
  \label{eq:bkg_outflow_rate}
\end{align}
$r_\text{in}^{X}$ is the fraction of the \emph{measured} region-$X$ count that
comes from migrated clusters. $r_\text{out}^{X}$ is the fraction of clusters
that \emph{should} have populated region $X$ but were lost to outward migration.

\subsection{Single-Interaction Results}
\label{subsec:bkg_flow_single}

Table~\ref{tab:bkg_flow_single} shows the integrated flow rates for jet20+jet30
(single interaction). All rates are below 3\%, with outflow and inflow of
comparable magnitude. The tight BDT regions (A, B) show outflow slightly larger
than inflow because the tight shower-shape cuts preferentially select clean
in-fiducial clusters that happen to migrate outward; the edge-crossed clusters
that migrate inward carry worse shower shapes and are rejected. Region C
(non-tight, iso) is the region most affected by inward migration.

\begin{table}[htbp]
  \centering
  \caption{Integrated inflow and outflow rates (jet20+jet30, single interaction)
    in the no-NPB classification.}
  \label{tab:bkg_flow_single}
  \begin{tabular}{lcc}
    \toprule
    Region & $\langle r_\text{in} \rangle$ & $\langle r_\text{out} \rangle$ \\
    \midrule
    A: tight + iso        & 1.01\% & 1.80\% \\
    B: tight + non-iso    & 0.91\% & 1.95\% \\
    C: non-tight + iso    & 2.70\% & 1.65\% \\
    D: non-tight + non-iso & 1.22\% & 1.33\% \\
    \bottomrule
  \end{tabular}
\end{table}

\subsection{Double-Interaction Results}
\label{subsec:bkg_flow_double}

The full-GEANT double-interaction sample \texttt{jet12\_double} shows
\emph{dramatically} larger migration rates: the vertex shift from the overlap
of two collisions pushes cluster $\eta$ by
$\Delta\eta \sim -\Delta z / (R_\text{EMCal} \cosh \eta) \sim 0.2$ on top of
the single-interaction cluster-position resolution. Integrated rates reach
several tens of percent in the tight-excluding non-tight regions (C, D), and
the inflow dominates strongly over the outflow:
the measured sideband populations are contaminated $\sim 30$--$40\%$ by
pileup-migrated clusters at high $\ET$. This is far larger than the
single-interaction effect and drives the bulk of the pileup-related $R$-factor
bias discussed below.

Note that \texttt{jet12\_double} represents the pure double-interaction limit.
At the nominal 1.5~mrad crossing angle, the cluster-weighted pileup fraction is
${\sim}7.9\%$ (from \texttt{calc\_pileup\_range.C}), so the effective bias in
data is the weighted average
$r^\text{eff} \approx 0.921 \, r^\text{single} + 0.079 \, r^\text{double}$.

\subsection{Impact on the ABCD $R$-Factor}
\label{subsec:bkg_rfactor}

The ABCD background estimate uses $R = N_B \cdot N_C / N_D$ to project sideband
counts into the signal region (with small signal-leakage corrections applied
at the next level). Because inflow is present in all three non-signal regions,
a substantial portion of the single-interaction bias cancels in the ratio.
Figure~\ref{fig:bkg_rfactor_bias} shows $R^\text{meas}/R^\text{native}$ where
$R^\text{meas}$ uses the measured sideband populations
$N_X^\text{meas} = N_X^\text{native} + N_X^\text{inflow}$ and $R^\text{native}$
uses $N_X^\text{native}$ only.

For single-interaction jet MC, the $R$-factor bias is $\lesssim 4\%$ across
most of the analysis range, rising to $\sim 10\%$ in the top two $\pT$ bins
where statistics are marginal. For pure double-interaction (\texttt{jet12\_double}),
the bias grows rapidly with $\ET$ and reaches $\times 2$--$\times 2.5$ at
$\ET \sim 20$--$25$~GeV. In the blended data scenario with 7.9\% pileup, the
effective bias is
$1.0 + 0.079 \times (R^\text{d}/R^\text{n} - 1)$, which translates to roughly
$\sim 6$--$12\%$ at mid-$\ET$ and potentially larger at high-$\ET$. This is
at the same order as PPG12's existing sideband systematic uncertainties and
warrants a follow-up systematic evaluation.

\subsection{Figures}
\label{subsec:bkg_flow_figs}

\begin{figure}[htbp]
  \centering
  \includegraphics[width=0.95\textwidth]{eta_migration_bkg_sideband_flow.pdf}
  \caption{Inflow and outflow rates (no-NPB classification) vs. reco $\ET$
    per ABCD region in jet MC. Single interaction (filled black / blue markers)
    produces rates $\lesssim 3\%$; the full-GEANT \texttt{jet12\_double}
    sample (open red / orange markers) drives much larger asymmetric
    inflow, reaching $\sim 30$--$70\%$ in regions C and D at high $\ET$.
    The asymmetry $r_\text{in} > r_\text{out}$ in the double-interaction case
    reflects the vertex-shift bias: pileup reconstructs clusters at a
    different $\eta$ than the leading truth particle.}
  \label{fig:bkg_sideband_flow}
\end{figure}

\begin{figure}[htbp]
  \centering
  \begin{minipage}[t]{0.49\textwidth}
    \centering
    \includegraphics[width=\textwidth]{eta_migration_bkg_net_flow.pdf}
    \caption{Net sideband bias $(N^\text{inflow} - N^\text{outflow}) /
      N^\text{native}$ vs. reco $\ET$ per region. Single interaction is
      consistent with zero at $\lesssim 2\%$; double interaction drives
      significant positive bias in regions C and D at high $\ET$.}
    \label{fig:bkg_net_flow}
  \end{minipage}
  \hfill
  \begin{minipage}[t]{0.49\textwidth}
    \centering
    \includegraphics[width=\textwidth]{eta_migration_bkg_rfactor_bias.pdf}
    \caption{Ratio $R^\text{meas}/R^\text{native}$ of the ABCD $R$-factor
      using measured ($N^\text{native} + N^\text{inflow}$) vs. native
      populations. Single interaction shows $< 10\%$ bias; pure double
      interaction (\texttt{jet12\_double}) shows a bias up to $\times 2.5$ at
      mid-$\ET$.}
    \label{fig:bkg_rfactor_bias}
  \end{minipage}
\end{figure}

\begin{figure}[htbp]
  \centering
  \begin{minipage}[t]{0.49\textwidth}
    \centering
    \includegraphics[width=\textwidth]{eta_migration_bkg_matrix.pdf}
    \caption{Migration matrix (reco $\eta$ vs.\ truth $\eta$ of the dominant
      contributor) for single-interaction jet MC. Narrow diagonal with small
      off-diagonal tails across the $\abseta = 0.7$ boundary (red dashed).
      Color axis is logarithmic.}
    \label{fig:bkg_migration_matrix}
  \end{minipage}
  \hfill
  \begin{minipage}[t]{0.49\textwidth}
    \centering
    \includegraphics[width=\textwidth]{eta_migration_bkg_matrix_double.pdf}
    \caption{Migration matrix for the full-GEANT double-interaction sample
      \texttt{jet12\_double}. The diagonal is visibly broader and extensive
      off-diagonal migration appears near and across the $\abseta = 0.7$
      boundary.}
    \label{fig:bkg_migration_matrix_double}
  \end{minipage}
\end{figure}

\subsection{Summary}
\label{subsec:bkg_flow_summary}

\begin{itemize}
  \item For single-interaction background MC the eta edge produces comparable
    inflow and outflow in every ABCD sideband ($\lesssim 3\%$), and the
    residual $R$-factor bias is $\lesssim 4\%$ across the analysis range.
  \item For pure double-interaction MC (\texttt{jet12\_double}) the inflow
    dominates strongly (tens of percent), and the $R$-factor bias reaches
    $\times 2$--$2.5$ at mid $\ET$.
  \item Applied with the measured 7.9\% pileup fraction at 1.5~mrad, the
    effective $R$-factor bias in data is at the $\sim 5$--$15\%$ level,
    comparable to other PPG12 ABCD systematic uncertainties and motivating
    a dedicated systematic evaluation.
\end{itemize}

\end{document}
